
\subsubsection{6.1 Key Finding: Propagation Without Lock-in}

Our results support a nuanced view of benchmark concentration in ML research. Citation networks strongly propagate benchmark standards---papers that cite each other share benchmarks at approximately 24 times the rate of random pairs (Cohen's d = 1.93). This confirms that evaluation practices spread through social and intellectual networks.

However, we find no evidence of the epistemic lock-in that critics fear. Concentration in year t-1 does not predict diversity decline in year t ($\beta$ = +1.60, p = 0.098). The positive coefficient, while not significant, suggests if anything that concentration is followed by diversification rather than further concentration.

\subsubsection{6.2 Theoretical Implications}

These findings favor a co-evolution model over a deterministic lock-in model. In the co-evolution view, benchmark concentration and research practices evolve together without causal feedback---concentration may rise and fall based on field dynamics (new benchmark releases, paradigm shifts) rather than self-reinforcing cycles.

The strong propagation mechanism (H-M3) combined with the absence of temporal feedback (H-M5) suggests that while the ML community transmits evaluation standards efficiently, it retains the capacity to shift these standards. The field is coordinated, not captured.

\subsubsection{6.3 Limitations}

\textbf{Small panel size.} With only 21 venue-years (18 after lag-drop), our panel regression has limited statistical power. The 95\% confidence interval for the temporal effect spans zero ([-0.37, 3.58]).

\textbf{Task labels as dataset proxy.} Papers With Code task labels (e.g., ''Image Classification'') are coarser than specific dataset names. This likely attenuates our concentration estimates and may mask dataset-level lock-in within task categories.

\textbf{Citation proxy.} We used task co-occurrence as a proxy for citation relationships due to API rate limits. While papers on the same task typically cite each other, this proxy conflates genuine intellectual influence with topical similarity.

\textbf{Observational design.} We document associations, not causal mechanisms. The temporal analysis provides stronger evidence than cross-sectional correlation, but cannot establish experimental causation.

\textbf{Venue scope.} Our analysis covers three general ML venues. Domain-specific venues (CVPR, ACL, etc.) may exhibit different dynamics.

\textbf{Coverage variability.} Papers With Code coverage varies substantially across venue-years, from 6 papers \citep{ICML2024} to 1,948 papers \citep{NeurIPS2021}. Low-coverage venue-years contribute noise to aggregated metrics.

\subsubsection{6.4 H-M4 Failure Analysis}

The failure of H-M4 (diversity-concentration correlation) warrants discussion. While ICML showed a strong significant negative correlation ($\rho$ = -0.90, p = 0.037), NeurIPS and ICLR did not show significant effects. Several explanations are possible:

\begin{enumerate}
\item Task label count may not capture actual evaluation diversity
\item 19 venue-years may be underpowered for detecting moderate effects
\item The true effect may be smaller than hypothesized
\item Venue-specific factors may moderate the concentration-diversity relationship
\end{enumerate}

